% Volume III: Attention and Publicity
% Part V. Publicity Without Comprehension
% Chapter 23: Visibility Is Not Knowledge
% Spine: proof-spines/ch023-spine.md

\chapter{Visibility Is Not Knowledge}
\label{ch:ch023}

The chapter begins by rejecting the identification
\[
\text{more recording}=\text{more knowledge}.
\]
A more adequate relation is
\[
K = f(V,I,C,A),
\]
\ToyModel\ where \(V\) is evidentiary volume, \(I\) indexing quality, \(C\) contextual integrity, and \(A\) available attention. The point is not that visibility is worthless, but that knowledge depends on the interaction of these variables. If \(V\) rises while \(I\), \(C\), and \(A\) are held fixed, the probability of locating and correctly interpreting a relevant item can fall rather than rise.

Visibility multiplies appearances without by itself settling cause, relevance, or evidentiary weight. A regime of disclosure can therefore leave observers overwhelmed, manipulable, or falsely confident. The volume opens by marking the distinction between what can be seen and what can be responsibly known, and by motivating the mediating work later chapters assign to indexing, attention-allocation, contextual boundaries, and organized collective inquiry.
